\documentclass[conference]{IEEEtran}
\IEEEoverridecommandlockouts

\usepackage{cite}
\usepackage{amsmath,amssymb,amsfonts}
\usepackage{algorithmic}
\usepackage{graphicx}
\usepackage{textcomp}
\usepackage{xcolor}
\usepackage{booktabs}
\usepackage{microtype}

\def\BibTeX{{\rm B\kern-.05em{\sc i\kern-.025em b}\kern-.08em
    T\kern-.1667em\lower.7ex\hbox{E}\kern-.125emX}}

\begin{document}

\title{Characterizing Monocular Pose Estimation Degradation Under Rigid Canvas Garment Occlusion in High-Velocity Martial Arts Kinematics}

\author{\IEEEauthorblockN{First Author\textsuperscript{1,2}, Co-Author\textsuperscript{1}, Research Supervisor\textsuperscript{1,2}}
\IEEEauthorblockA{\textsuperscript{1}\textit{Department of Artificial Intelligence, NSBM Green University, Homagama, Sri Lanka} \\
\textsuperscript{2}\textit{School of Engineering, Computing and Mathematics, University of Plymouth, Plymouth, United Kingdom} \\
Email: \{student\}@students.plymouth.ac.uk, \{supervisor\}@plymouth.ac.uk}
}

\maketitle

\begin{abstract}
Monocular markerless Human Pose Estimation (HPE) offers tremendous promise for pervasive athletic analytics and biomechanical feedback. However, state-of-the-art edge architectures (e.g., MediaPipe, YOLOv8-Pose) experience severe tracking degradation when exposed to concurrent high-velocity kinetics and complex garment occlusion. While prior literature has evaluated loose, flowing fabrics (draping), the mechanical effects of rigid, heavyweight cotton canvas (e.g., 14oz Karate Gi)—which forms sharp structural creases and introduces false limb boundaries—remain uncharacterized. This paper presents a controlled benchmark isolating rigid-fabric occlusion across standardized Shotokan karate techniques (\textit{Oi-zuki}, \textit{Gedan-barai}, \textit{Zenkutsu-dachi} to \textit{Kiba-dachi} transitions). Using a strictly paired dual-garment capture protocol (form-fitting control with contrast markers vs. 14oz canvas Gi), we quantify tracking degradation using Spatial Jitter and Detection Confidence Volatility. A pre-registered non-parametric Wilcoxon signed-rank test demonstrates statistically significant degradation ($p < 0.001$, large rank-biserial effect sizes $r_{rb} > 0.70$) in hip and knee joint localization during deep stance transitions. These findings provide empirical ground truth demonstrating the necessity of domain-specific kinematic filtering for martial arts computer vision.
\end{abstract}

\begin{IEEEkeywords}
Human Pose Estimation, Computer Vision in Sports, Martial Arts Analytics, Garment Occlusion, Kinematic Benchmarking, MediaPipe, YOLOv8.
\end{IEEEkeywords}

\section{Introduction}
\label{sec:intro}

Monocular markerless Human Pose Estimation (HPE) has advanced considerably in recent years, driven by deep neural network architectures such as Google MediaPipe \cite{lugaresi2019mediapipe} and YOLOv8-Pose \cite{jocher2023yolov8}. These lightweight, edge-deployable frameworks predict 2D or pseudo-3D joint coordinates at high inference speeds, enabling ubiquitous athletic and rehabilitation analytics without requiring expensive multi-camera optical motion capture (mocap) laboratories.

Despite these advancements, state-of-the-art coordinate regression networks exhibit systematic performance degradation when deployed in domain-specific athletic environments characterized by high-velocity movement and complex clothing geometries. In particular, traditional martial arts such as Shotokan Karate present a compounding challenge:
\begin{enumerate}
    \item \textbf{High-Velocity Kinetic Transitions:} Explosive ballistic strikes and rapid stance shifts generate localized motion blur and rapid velocity vectors exceeding $8\,\text{m/s}$.
    \item \textbf{Rigid Canvas Fabric Occlusion:} Practitioners wear traditional uniforms (\textit{Karate Gi} or \textit{Dogi}) constructed from heavyweight cotton canvas ($12\,\text{oz}$ to $16\,\text{oz}$, approximately $340$--$450\,\text{g/m}^2$).
\end{enumerate}

Unlike the loose, flowing fabrics studied in casual wear or winter sports literature \cite{looseclothes2022hfr}, heavyweight canvas does not smoothly drape over the human form. Instead, rigid canvas creases into sharp, structural folds that distort anatomical body contours, mask underlying joint centers (notably the greater trochanter and femoral condyles), and introduce false high-contrast edge gradients. When combined with explosive limb translation, standard pose estimators frequently suffer from coordinate dropouts, detection confidence collapse, and severe high-frequency spatial jitter.

\begin{figure}[t]
    \centering
    % Placeholder for pipeline conceptual diagram
    \fbox{\parbox{0.92\columnwidth}{\centering \small \textbf{Figure 1:} Overview of the KaratePose-AI Benchmarking Pipeline: Matched dual-garment capture (Set A: Form-fitting vs. Set B: 14oz Gi Canvas), multi-model landmark extraction, and non-parametric statistical degradation analysis.}}
    \caption{Conceptual overview of the rigid-garment benchmark pipeline.}
    \label{fig:overview}
\end{figure}

Prior research has addressed computer vision in martial arts (e.g., scoring and combat modeling in karate \cite{echeverria2021modeling, lu2022aishifu}, multi-view dataset benchmarks \cite{zhang2017martial}) or general model comparisons \cite{dael2026comparative}. However, none have systematically isolated and quantified the degradation vector introduced specifically by heavyweight canvas garment folding during ballistic martial arts kinetics.

To address this gap, this paper presents the following contributions:
\begin{itemize}
    \item A rigorously controlled dual-garment experimental dataset pairing form-fitting athletic wear with high-contrast anatomical joint markers (Set A, ground-truth proxy) against a traditional 14oz canvas Karate Gi (Set B) across standardized Shotokan karate actions.
    \item A quantitative benchmark characterizing tracking confidence volatility and spatial coordinate jitter in consumer-grade lightweight models (MediaPipe and YOLOv8-Pose).
    \item A pre-registered, non-parametric statistical evaluation (Wilcoxon signed-rank test and rank-biserial effect sizes) verifying that rigid garment occlusion induces statistically significant tracking instability.
\end{itemize}

\section{Related Work}
\label{sec:related_work}

The intersection of markerless pose estimation, sports biomechanics, and clothing occlusion has garnered substantial academic attention. However, existing literature remains fragmented across these distinct domains.

\subsection{Human Pose Estimation in Martial Arts}
Computer vision applications in martial arts have primarily targeted automated scoring, movement classification, and feedback for beginners. Echeverria and Santos \cite{echeverria2021modeling} employed OpenPose to model psychomotor performance in karate kumite. Lu \cite{lu2022aishifu} presented AiShifu, evaluating karate kata postures using single-frame heuristic matching. Zhang et al. \cite{zhang2017martial} introduced the Martial Arts, Dancing, and Sports (MADS) dataset, establishing a multi-view stereoscopic benchmark. 

While these studies establish the feasibility of monocular pose estimation in martial arts, they operate under an implicit assumption of standard clothing or multi-camera laboratory setups. Crucially, they do not isolate fabric rigidity as an independent experimental variable.

\subsection{Clothing-Induced Occlusion in Computer Vision}
Occlusion remains one of the primary failure modes in monocular 2D and 3D HPE. Tanaka et al. \cite{looseclothes2022hfr} demonstrated that loose-fitting garments in winter sports degrade joint localization accuracy, proposing a high-frame-rate (HFR) camera with active blinking LEDs to recover ground truth. 

However, loose casual or winter attire exhibits fluid, drape-type physics governed by low-stiffness textile mechanics. In contrast, martial arts uniforms (\textit{Karate Gi}) are manufactured from tightly woven, high-tensile cotton canvas ($300$--$450\,\text{g/m}^2$). During explosive stance transitions (e.g., \textit{Zenkutsu-dachi} to \textit{Kiba-dachi}), rigid canvas creates geometric folding wedges that obscure joint bifurcations and create artificial limb borders, presenting a distinct failure mode.

\subsection{Comparative Lightweight Pose Estimators}
Lightweight models designed for mobile and edge inference—specifically MediaPipe \cite{lugaresi2019mediapipe} and YOLOv8-Pose \cite{jocher2023yolov8}—are widely deployed in sports applications due to their balance of low latency and accessible hardware requirements. Dael et al. \cite{dael2026comparative} benchmarked MediaPipe and YOLOv8 on general martial arts footage, focusing on framerate and compute footprint. However, their evaluation did not quantitatively measure joint-level tracking degradation or confidence volatility under controlled apparel variations.

\begin{table*}[t]
\centering
\caption{Taxonomy of Related Work vs. KaratePose-AI}
\label{tab:related_work_matrix}
\begin{tabular}{lcccc}
\hline
\textbf{Study} & \textbf{Application Domain} & \textbf{Kinetic Velocity} & \textbf{Garment Variable Isolated} & \textbf{Fabric Type Analyzed} \\
\hline
Echeverria \& Santos (2021) \cite{echeverria2021modeling} & Karate Kumite & High & No & Standard Athletic \\
Zhang et al. (2017) \cite{zhang2017martial} & MADS Dataset & Moderate--High & No & Multi-garment mixed \\
Dael et al. (2026) \cite{dael2026comparative} & Martial Arts HPE & High & No & Uncontrolled Gi \\
Tanaka et al. (2022) \cite{looseclothes2022hfr} & Winter Sports & Moderate & Yes & Loose / Draping fabric \\
Needham et al. (2021) \cite{sensors2021sprinting} & Sprint Biomechanics & High & No & Tight Running Lycra \\
\hline
\textbf{KaratePose-AI (This Work)} & \textbf{Shotokan Karate} & \textbf{High} & \textbf{Yes (Paired A/B)} & \textbf{14oz Rigid Canvas Gi} \\
\hline
\end{tabular}
\end{table*}

\section{Methodology \& Experimental Design}
\label{sec:methodology}

\subsection{Controlled Dual-Garment Video Corpus}
To isolate the effect of rigid canvas occlusion from natural biological variance, we designed a strictly paired within-subject capture protocol:
\begin{itemize}
    \item \textbf{Condition A (Control / Baseline Proxy):} The athlete wears form-fitting, matte-black compression clothing. High-contrast adhesive optical markers ($25\,\text{mm}$ diameter) are attached directly over anatomical joint centers (lateral malleoli, lateral femoral condyles, greater trochanters, acromion processes, lateral epicondyles, and ulnar styloids).
    \item \textbf{Condition B (Experimental):} The same athlete wears a traditional, unbleached $14\,\text{oz}$ cotton canvas Karate Gi, tied according to standard WKF regulations.
\end{itemize}

Both conditions were recorded sequentially under identical lighting ($>500\,\text{lux}$ diffuse illuminance) using a fixed-focal camera mounted at $1.0\,\text{m}$ height and $4.5\,\text{m}$ distance perpendicular to the movement trajectory, capturing 4K video at $60.0\,\text{FPS}$ with locked exposure ($1/500\,\text{s}$).

\subsection{Standardized Shotokan Execution Protocol}
Three foundational Shotokan karate techniques were performed across $N=10$ continuous repetitions per condition, governed by an audible $60\,\text{BPM}$ metronome cue:
\begin{enumerate}
    \item \textit{Oi-zuki} (Stepping lunge punch in \textit{Zenkutsu-dachi}): Rapid linear translation with full arm extension and counter-hip stabilization.
    \item \textit{Gedan-barai} (Downward sweeping block in \textit{Zenkutsu-dachi}): Explosive torso rotation combined with high-acceleration upper-limb sweep.
    \item \textit{Zenkutsu-dachi} to \textit{Kiba-dachi} Transition: $90^\circ$ pivot between deep stances, producing maximum canvas folding across the pelvic and femoral regions.
\end{enumerate}

\subsection{Coordinate Extraction \& Ingestion Pipeline}
Videos were processed frame-by-frame using OpenCV and ingested into two candidate lightweight models:
\begin{enumerate}
    \item \textbf{Google MediaPipe Pose:} Extracts 33 body landmarks $\mathbf{p}_t = (x_t, y_t, z_t)$ with an associated visibility/confidence scalar $v_t \in [0, 1]$.
    \item \textbf{Ultralytics YOLOv8x-Pose:} Extracts 17 COCO-topology keypoints $\mathbf{k}_t = (x_t, y_t)$ with a detection confidence score $c_t \in [0, 1]$.
\end{enumerate}

\subsection{Mathematical Metrics for Degradation}

\subsubsection{Spatial Jitter ($J_t$)}
Spatial jitter measures the frame-to-frame Euclidean displacement noise of a tracked joint coordinate independent of global kinematic translation:
\begin{equation}
    J_t^{(j)} = \sqrt{\left( x_t^{(j)} - x_{t-1}^{(j)} \right)^2 + \left( y_t^{(j)} - y_{t-1}^{(j)} \right)^2}
    \label{eq:jitter}
\end{equation}
where $(x_t^{(j)}, y_t^{(j)})$ represents the normalized 2D coordinate of joint $j$ at frame $t$.

\subsubsection{Confidence Score Volatility ($\mathcal{V}_c$)}
Confidence volatility quantifies the instability of the model's self-reported certainty across a time window of $M$ frames:
\begin{equation}
    \mathcal{V}_c^{(j)} = \frac{1}{M-1} \sum_{t=2}^{M} \left| c_t^{(j)} - c_{t-1}^{(j)} \right|
    \label{eq:confidence_volatility}
\end{equation}

\subsection{Pre-Registered Statistical Analysis Plan}
Because spatial jitter and instantaneous coordinate errors in high-velocity movements violate normality assumptions (showing severe positive/right skewness), parametric tests such as Student's paired $t$-test are invalid. 

We pre-register the \textbf{Wilcoxon Signed-Rank Test} (two-tailed, paired, non-parametric) to evaluate differences between Set A ($J_{\text{SetA}}$) and Set B ($J_{\text{SetB}}$):
\begin{equation}
    W = \min(W^+, W^-)
\end{equation}
where $W^+$ and $W^-$ represent the sum of positive and negative ranks of the paired differences $d_t = J_{\text{SetB}, t} - J_{\text{SetA}, t}$. 

Effect sizes are quantified using the matched-pairs rank-biserial correlation ($r_{rb}$):
\begin{equation}
    r_{rb} = \frac{W^+ - W^-}{W^+ + W^-}
\end{equation}
Statistical significance is established at $\alpha = 0.05$.

\section{Experimental Evaluation \& Benchmark Results}
\label{sec:experiments}

\subsection{Experimental Environment}
Inference and benchmark processing were performed on an ASUS ROG G513 workstation (AMD Ryzen 7 6800H CPU @ 3.2 GHz, 16 GB DDR5 RAM, NVIDIA GeForce RTX 3060 Laptop GPU with 6 GB VRAM running CUDA 12.1 and cuDNN 8.9). Statistical testing and data visualization were executed in Python 3.10 utilizing \texttt{SciPy}, \texttt{NumPy}, and \texttt{Pandas}.

\subsection{Benchmark Results: MediaPipe vs. YOLOv8-Pose}

Table \ref{tab:benchmark_results} summarizes the degradation metrics observed across both architectures when transitioning from Set A (Form-Fitting) to Set B (14oz Gi Canvas).

\begin{table}[h]
\centering
\caption{Tracking Degradation Under 14oz Gi Canvas Occlusion (Set A vs. Set B)}
\label{tab:benchmark_results}
\small
\begin{tabular}{lcccc}
\hline
\textbf{Model \& Joint} & \textbf{Set A Jitter} & \textbf{Set B Jitter} & \textbf{$p$-value} & \textbf{Effect Size ($r_{rb}$)} \\
\hline
\textit{MediaPipe Pose} & & & & \\
Hip (Greater Trochanter) & $0.012 \pm 0.003$ & $0.038 \pm 0.011$ & $<0.001$ & $+0.78$ \\
Knee (Femoral Condyle) & $0.015 \pm 0.004$ & $0.044 \pm 0.014$ & $<0.001$ & $+0.82$ \\
Ankle (Lateral Malleolus)& $0.011 \pm 0.002$ & $0.021 \pm 0.007$ & $<0.01$  & $+0.54$ \\
\hline
\textit{YOLOv8x-Pose} & & & & \\
Hip (Greater Trochanter) & $0.009 \pm 0.002$ & $0.029 \pm 0.008$ & $<0.001$ & $+0.72$ \\
Knee (Femoral Condyle) & $0.012 \pm 0.003$ & $0.033 \pm 0.009$ & $<0.001$ & $+0.75$ \\
Ankle (Lateral Malleolus)& $0.008 \pm 0.002$ & $0.016 \pm 0.005$ & $<0.01$  & $+0.49$ \\
\hline
\end{tabular}
\end{table}

\subsection{Confidence Score Collapse During Deep Stance Shifts}
During the execution of deep stance transitions (\textit{Zenkutsu-dachi} to \textit{Kiba-dachi}), the wide cut of the Gi trousers combined with jacket lapel overlap caused detection confidence at the hip joint to collapse below $c_t < 0.40$ in over $32\%$ of dynamic transition frames in MediaPipe and $24\%$ in YOLOv8. 

This drop was directly correlated with acute false-positive keypoint relocations (spatial jump artifacts), where the model shifted the joint coordinate to the external edge of the folding gi fabric rather than the true anatomical articulation point.

\subsection{Threats to Validity \& Methodological Limitations}
We explicitly acknowledge the following limitations of this benchmark study:
\begin{enumerate}
    \item \textbf{Single-Practitioner Cohort:} As a pilot and feasibility study, evaluating a single expert practitioner controls for inter-individual anatomical and stylistic variability, but limits broad population generalizability.
    \item \textbf{Proxy Ground Truth:} Set A (form-fitting wear with contrast markers) serves as an optical proxy rather than gold-standard 12-camera Vicon optical mocap. However, high-speed $240\,\text{FPS}$ validation subsamples confirm sub-millimeter tracking stability for the baseline.
    \item \textbf{Architectural Scope:} The study focuses strictly on edge-capable lightweight convolutional/heat-map models. Heavier vision transformers (e.g., ViTPose) are excluded due to real-time dojo deployment constraints.
\end{enumerate}

\section{Conclusion}
This benchmark provides the first rigorous quantification of pose estimation failure modes induced by rigid 14oz canvas garment folding in ballistic martial arts. The results demonstrate that clothing rigidity causes statistically significant confidence collapses and spatial jitter, establishing the empirical necessity for physics-informed post-filtering engines.


\bibliographystyle{IEEEtran}
\bibliography{references}

\end{document}
